% !TeX program = xelatex
% Compile twice with XeLaTeX.
\documentclass[UTF8,11pt,a4paper,fontset=fandol]{ctexart}
\usepackage[margin=26mm]{geometry}
\usepackage{amsmath,amssymb,amsthm,mathtools}
\usepackage{xcolor}
\usepackage[colorlinks=true,linkcolor=blue!60!black,citecolor=blue!60!black,urlcolor=blue!60!black]{hyperref}
\setlength{\emergencystretch}{2em}
\newtheorem{theorem}{定理}
\newtheorem{corollary}[theorem]{推论}
\theoremstyle{remark}
\newtheorem{remark}[theorem]{说明}
\newtheorem{example}[theorem]{例}
\renewcommand{\proofname}{证明}
\newcommand{\PP}{\mathcal P}
\newcommand{\CC}{\mathcal C}
\newcommand{\HH}{\mathcal H}
\newcommand{\Tr}{\operatorname{Tr}}
\newcommand{\rank}{\operatorname{rank}}
\newcommand{\norm}[1]{\left\lVert #1\right\rVert}
\title{三角投影交换子作为 Hilbert--Schmidt 空间上算子的 Schatten 类刻画}
\author{}
\date{}

\begin{document}
\maketitle
\vspace{-2em}

设 $H$ 为可分复 Hilbert 空间，$P_0=0<P_1<P_2<\cdots$ 为有限秩正交投影，
并且 $P_n\to I$ 强算子收敛。记
\[
 E_n=P_n-P_{n-1},\qquad d_n=\rank E_n,\qquad \HH=S_2(H).
\]
在 Hilbert 空间 $\HH$ 上，令 $\PP$ 为块上三角正交投影：
\[
 \PP(A)=\sum_{i\leq j}E_iAE_j.
\]
给定 $X\in B(H)$，记 $L_X(A)=XA$，并考虑
\[
 \CC_X=[L_X,\PP],\qquad
 \CC_X(A)=X\PP(A)-\PP(XA).
\]
\textbf{本文的 $S_p(\HH)=S_p(S_2(H))$ 指算子 $\CC_X$ 本身的 Schatten 类，}
其定义使用 Hilbert 空间 $\HH$ 上的奇异值；它与 $\CC_X:S_p(H)\to S_p(H)$
是否有界是两个不同问题。

\section*{充要条件和精确范数公式}

\begin{theorem}\label{thm:main}
对任意 $0<p<\infty$（特别是 $1<p<\infty$），下列条件等价：
\[
 \boxed{\quad
 \CC_X\in S_p\bigl(S_2(H)\bigr)
 \quad\Longleftrightarrow\quad
 \sum_{n=1}^{\infty}d_n\norm{[X,P_n]}_{S_p(H)}^p<\infty.
 \quad}
\]
而且有精确等式
\begin{equation}\label{eq:main}
 \norm{\CC_X}_{S_p(S_2(H))}^p
 =\sum_{n=1}^{\infty}d_n\norm{[X,P_n]}_{S_p(H)}^p.
\end{equation}
等价地，该条件可以写为
\begin{equation}\label{eq:corners}
 \boxed{\quad
 \sum_{n=1}^{\infty}d_n\left(
 \norm{(I-P_n)XP_n}_{S_p(H)}^p
 +\norm{P_nX(I-P_n)}_{S_p(H)}^p\right)<\infty.
 \quad}
\end{equation}
\eqref{eq:corners} 中的和也精确等于 \eqref{eq:main} 的左侧。
对于 $0<p<1$，这里使用 Schatten 拟范数。
\end{theorem}

\begin{proof}
\textbf{第一步：将交换子分解为角块交换子的直和。}
对每个 $n$，取 $E_nH$ 的一组标准正交基
$\{\eta_{n,r}:1\leq r\leq d_n\}$。
所有这些向量合起来构成 $H$ 的标准正交基。
由 Hilbert--Schmidt 范数的列表示，映射
\[
 U:\HH\longrightarrow\bigoplus_{n\geq1}\bigoplus_{r=1}^{d_n}H,
 \qquad U(A)=(A\eta_{n,r})_{n,r}
\]
是酉算子，因为
\[
 \norm A_2^2=\sum_{n\geq1}\sum_{r=1}^{d_n}\norm{A\eta_{n,r}}_H^2,
\]
且任意平方可和的向量列均唯一确定一个 Hilbert--Schmidt 算子。

固定 $n,r$。由于 $\eta_{n,r}\in E_nH$，有
\[
 \PP(A)\eta_{n,r}=P_nA\eta_{n,r}.
\]
因此
\begin{align*}
 \CC_X(A)\eta_{n,r}
 &=XP_nA\eta_{n,r}-P_nXA\eta_{n,r}\\
 &=[X,P_n]A\eta_{n,r}.
\end{align*}
于是得到了真正的酉等价
\begin{equation}\label{eq:unitary}
 \boxed{\quad
 U\CC_XU^{-1}
 =\bigoplus_{n\geq1}\bigoplus_{r=1}^{d_n}[X,P_n].
 \quad}
\end{equation}
换言之，$[X,P_n]$ 的每个非零奇异值，在这一分量中恰好重复 $d_n$ 次。

\medskip
\textbf{第二步：对奇异值求 $p$ 次幂和。}
每个 $K_n=[X,P_n]$ 都是有限秩算子，因为 $P_n$ 有限秩且 $X$ 有界。
直和算子的奇异值由各分量的非零奇异值汇集而成，只要这些奇异值趋于零。
若 $\sum_nd_n\norm{K_n}_p^p<\infty$，则
\[
 \norm{K_n}^p\leq\norm{K_n}_p^p\longrightarrow0.
\]
因此 \eqref{eq:unitary} 的直和是紧算子，其奇异值 $p$ 次幂之和正是
$\sum_nd_n\norm{K_n}_p^p$，证明了充分性及公式 \eqref{eq:main}。
反过来，若 $\CC_X\in S_p(\HH)$，
则其有限分量直和的奇异值 $p$ 次幂和一致有界；
对分量数取极限，即得同一公式和该级数的收敛性。
这一论证对所有 $p>0$ 都成立。

\medskip
\textbf{第三步：改写为两个方向的角块条件。}
设
\[
 B_n=(I-P_n)XP_n,
 \qquad A_n=P_nX(I-P_n).
\]
相对于 $H=P_nH\oplus(I-P_n)H$，有
\[
 K_n=[X,P_n]=\begin{pmatrix}0&-A_n\\B_n&0\end{pmatrix},
 \qquad
 K_n^*K_n=\begin{pmatrix}B_n^*B_n&0\\0&A_n^*A_n\end{pmatrix}.
\]
从而对任意 $p>0$，
\[
 \norm{K_n}_p^p=\norm{B_n}_p^p+\norm{A_n}_p^p.
\]
代入 \eqref{eq:main}，即得 \eqref{eq:corners}。
\end{proof}

\begin{remark}[为何会出现重数 $d_n$]
在 $S_2(H)E_n$ 上，左乘 $[X,P_n]$ 对 $E_nH$ 的每一个基向量所对应的列，
都独立作用一次。因此一个块列包含 $d_n$ 个相同的 Hilbert 空间算子分量。
重数为一时 $d_n=1$；一般的有限维块情形不能省略这个因子。
\end{remark}

\begin{corollary}[自伴符号]
若 $X=X^*$，则两个角块互为伴随，因而
\[
 \norm{\CC_X}_{S_p(S_2(H))}^p
 =2\sum_{n\geq1}d_n\norm{(I-P_n)XP_n}_{S_p(H)}^p.
\]
所以自伴情形只需检验一个方向的角块的加权 $p$ 次幂可和性。
\end{corollary}

\section*{重数为一的矩阵形式及 Hilbert--Schmidt 特例}

在标准巢 $H=\ell_2(\mathbb N)$、
$P_n$ 投影到 $\operatorname{span}\{e_1,\ldots,e_n\}$ 的情形，
令 $X=(x_{ij})$，则 $d_n=1$，上述条件成为
\begin{equation}\label{eq:matrix}
 \sum_{n=1}^{\infty}\left(
  \norm{(x_{ij})_{i>n,\,j\leq n}}_{S_p}^p
  +\norm{(x_{ij})_{i\leq n,\,j>n}}_{S_p}^p
 \right)<\infty.
\end{equation}
这里的矩阵表示相应两个 Hilbert 空间之间的角块算子，
范数为其奇异值的 Schatten 范数。

当 $p=2$ 时可以把这一条件完全展开：
\begin{align}
 \norm{\CC_X}_{S_2(S_2(H))}^2
 &=\sum_{n\geq1}\left(
   \sum_{i>n,\,j\leq n}|x_{ij}|^2
   +\sum_{i\leq n,\,j>n}|x_{ij}|^2\right)\notag\\
 &=\boxed{\ \sum_{i\ne j}|i-j|\,|x_{ij}|^2\ }.\label{eq:hs}
\end{align}
事实上，当 $i>j$ 时，$x_{ij}$ 恰好出现在 $n=j,j+1,\ldots,i-1$
这 $i-j$ 个下角块中；当 $i<j$ 时，计数同理。
所以
\[
 \CC_X\in S_2(S_2(H))
 \quad\Longleftrightarrow\quad
 \sum_{i\ne j}|i-j|\,|x_{ij}|^2<\infty.
\]
对于 $p\ne2$，应使用 \eqref{eq:matrix} 中的角块奇异值范数；
上述计算不提供将其改成逐项 $|x_{ij}|^p$ 加权和的依据。

更一般地，在有限维块情形，令 $r_n=\rank P_n$、$r_0=0$，则
\begin{equation}\label{eq:blockhs}
 \norm{\CC_X}_{S_2(S_2(H))}^2
 =\sum_{i\ne j}|r_{i-1}-r_{j-1}|\,
   \norm{E_iXE_j}_{S_2(H)}^2.
\end{equation}
这是因为每个 $E_iXE_j$ 的权重为
$\sum_{n=\min(i,j)}^{\max(i,j)-1}d_n$。

\section*{与紧性、符号的 Schatten 条件的区别}

由 \eqref{eq:unitary}，还立即得到
\[
 \CC_X\text{ 在 }S_2(H)\text{ 上紧}
 \quad\Longleftrightarrow\quad
 \norm{[X,P_n]}_{B(H)}\longrightarrow0.
\]
这是原文定理 2.1(ii) 的紧性条件的角块表述。
本题要求的 Schatten 条件则进一步要求角块的\emph{全部奇异值}满足
\eqref{eq:main} 的加权可和性。
本文的直和分解与 \cite[p.~279, Remark (2)]{Power} 对 Hankel 分量的
分解相呼应，但 \eqref{eq:unitary} 是直接针对整个交换子得到的公式。

\begin{example}[即使 $X$ 为秩一，也不保证交换子属于给定 Schatten 类]
固定 $p>0$，在标准重数为一巢上令
\[
 u_i=i^{-1/2-1/p},\qquad
 X=u\otimes e_1^*.
\]
由于 $u\in\ell_2$，$X$ 是有界秩一算子，属于每个 $S_q(H)$，$q>0$。
然而对 $n\geq1$，有
\[
 P_nX(I-P_n)=0,\qquad
 (I-P_n)XP_n=((I-P_n)u)\otimes e_1^*.
\]
所以
\[
 \norm{[X,P_n]}_{S_p(H)}^p
 =\left(\sum_{i>n}i^{-1-2/p}\right)^{p/2}
 \asymp\frac1n.
\]
由调和级数发散及定理 \ref{thm:main}，
$\CC_X\notin S_p(S_2(H))$。
同时角块算子范数趋于零，故这个交换子仍然是紧算子。
这说明条件不能简化为 $X\in S_p(H)$，甚至 $X$ 为有限秩也不够。
\end{example}

\begin{remark}[对角部分不影响结论]
若 $D\in B(H)$ 为相对于 $(E_n)$ 的块对角算子，则
$[D,P_n]=0$，因此 $\CC_{X+D}=\CC_X$。
本题的条件和范数均只取决于 $X$ 的非对角块。
\end{remark}

\begin{remark}[只假定交换子有界时仍然成立]
若沿用前一个问题中允许形式块矩阵符号的设置，只假定
$\CC_X$ 已延拓为 $S_2(H)$ 上的有界算子，定理仍然成立。
事实上，以 $\xi\otimes\eta_{n,r}^*$ 作输入，逐列公式表明形式矩阵
$[X,P_n]$ 唯一确定一个范数不超过 $\norm{\CC_X}$ 的有界算子。
它相对于 $P_nH\oplus(I-P_n)H$ 只有两个非对角角块，故仍然为有限秩。
于是 \eqref{eq:unitary} 及其后的奇异值计算完全适用。
因此公式要求的是角块算子的有界性及可和性，而不依赖于对角符号是否有界。
\end{remark}

\begin{thebibliography}{9}
\bibitem{Power}
S. Power,
\emph{Commutators with the triangular projection and Hankel forms on nest algebras},
J. London Math. Soc. (2) \textbf{32} (1985), 272--282.
\href{https://doi.org/10.1112/jlms/s2-32.2.272}{doi:10.1112/jlms/s2-32.2.272}.
\end{thebibliography}

\end{document}
